\subsection{卷积与转置}

我们回到 No.~1 开头处的假设与记号，但此外假设 \(\beta\) 是以 \(\chi\) 为乘子的相对不变测度；于是 \(\chi\) 是 G 上的连续函数。

命题 7. --- 设 \(f\) 是局部 \(\beta\)-可积的函数，定义在 \(X\) 上，\(\nu\) 是 \(X\) 上的测度，\(\mu\) 是 \(G\) 上的测度。假设：

\begin{enumerate}
\def\labelenumi{(\roman{enumi})}
\item
  \(\mu\) 与 \(f\) 可作卷积，且 No.~1 的公式 (3) 在局部 \(\beta\)-几乎处处定义了一个卷积 \(\mu *^{\beta} f\)。
\item
  \(\chi \cdot \check{\mu}\) 与 \(\nu\) 可作卷积。
\item
  函数 \(g(s,x) = f(s^{-1}x)\chi (s^{-1})\) 是 \((\mu \otimes \nu)\)-可积的。
\end{enumerate}

则 \(f\) 对 \((\chi \cdot \breve{\mu}) * \nu\) 是本质可积的，由 (3) 定义的函数 \(\mu *^{\beta} f\) 是 \(\nu\)-可积的，且

\[
\nu (\mu * ^ {\beta} f) = \big ((\chi \cdot \check {\mu}) * \nu \big) (f).\tag{7}
\]

由于 \(g(s, x)\) 对 \(\mu \otimes \nu\) 可积，函数 \(f(sx)\) 对 \((\chi \cdot \breve{\mu}) \otimes \nu\) 是本质可积的，且 \(f\) 对 \((\chi \cdot \breve{\mu}) * \nu\) 是本质可积的。由 Lebesgue--Fubini 定理，\(\mu *^{\beta} f = \int g(s, x) d\mu(s)\) 是 \(\nu\)-可积的，且

\[
\begin{array}{l} \nu (\mu * ^ {\beta} f) = \iint f (s ^ {- 1} x) \chi (s ^ {- 1}) d \mu (s) d \nu (x) \\ = \iint f (s x) \chi (s) d \stackrel {\vee} {\mu} (s) d \nu (x) = \big ((\chi \cdot \stackrel {\vee} {\mu}) * \nu \big) (f). \end{array}
\]

例. --- 1) 由命题 3 以及 §1, No.~4 命题 5 的推论，可取 \(f \in \mathcal{C}(\mathrm{X})\)，\(\nu \in \mathcal{C}'(\mathrm{X})\) 与 \(\mu \in \mathcal{C}'(\mathrm{G})\)。这时公式 (7) 表示自同态 \(\nu \mapsto (\chi \cdot \breve{\mu}) * \nu\)（\(\mathcal{C}'(\mathrm{X})\) 的）是自同态 \(f \mapsto \mu * f\)（\(\mathcal{C}(\mathrm{X})\) 的）的转置。

\begin{enumerate}
\def\labelenumi{\arabic{enumi})}
\setcounter{enumi}{1}
\item
  可取 \(f \in \mathcal{K}(\mathrm{X})\)，\(\nu \in \mathcal{M}(\mathrm{X})\) 与 \(\mu \in \mathcal{C}'(\mathrm{G})\)（依据命题 3、§3, No.~2 的命题 8，以及连续函数 \(g(s,x)\) 的支集与 \(\mu \otimes \nu\) 的支集相交于一个紧集这一备注）。这时公式 (7) 表示自同态 \(\nu \mapsto (\chi \cdot \breve{\mu}) * \nu\)（\(\mathcal{M}(\mathrm{X})\) 的）是自同态 \(f \mapsto \mu * f\)（\(\mathcal{K}(\mathrm{X})\) 的）的转置。
\item
  若 G 真作用于 X，则可取 \(f \in \mathcal{K}(X)\)，\(\nu \in \mathcal{C}'(X)\) 与 \(\mu \in \mathcal{M}(G)\)（依据命题 4、§3, No.~2 的命题 8 以及例 2 中的同样备注）。
\end{enumerate}

命题 8. --- 设 \(f\) 与 \(g\) 是两个局部 \(\beta\)-可积的函数，定义在 \(X\) 上，并设 \(\mu \in \mathcal{M}(G)\)。假设：

\begin{enumerate}
\def\labelenumi{(\roman{enumi})}
\item
  \(\mu\) 与 \(f\) 可作卷积，且 No.~1 的公式 (3) 在局部 \(\beta\)-几乎处处定义了一个卷积 \(\mu *^{\beta} f\)。
\item
  \(\chi \cdot \breve{\mu}\) 与 \(g\) 可作卷积，且 No.~1 的公式 (3)（把 \(\mu\) 换成 \(\chi \cdot \breve{\mu}\)、把 \(f\) 换成 \(g\)）在局部 \(\beta\)-几乎处处定义了一个卷积 \((\chi \cdot \breve{\mu}) *^{\beta} g\)。
\item
  存在一个函数 \(\psi\)（定义在 \(\mathbf{G}\) 上），它局部 \(\mu\)-几乎处处等于 1，使得函数
\end{enumerate}

\[
h (s, x) = g (x) f \left(s ^ {- 1} x\right) \chi \left(s ^ {- 1}\right) \psi (s)
\]

是 \((\mu \otimes \beta)\)-可积的。

则函数 \(g(x)\big((\mu * ^ {\beta} f)(x)\big)\) 与 \(f(x)\big(\big((\chi \cdot \breve{\mu}) * ^ {\beta} g\big)(x)\big)\) 是本质 \(\beta\)-可积的，且

\[
\int f (x) \big (\big ((\chi \cdot \breve {\mu}) * ^ {\beta} g \big) (x) \big) d \beta (x) = \int g (x) \big ((\mu * ^ {\beta} f) (x) \big) d \beta (x).\tag{8}
\]

因为，由 (iii) 与 Lebesgue--Fubini 定理，函数

\[
x \mapsto g (x) \int f (s ^ {- 1} x) \chi (s ^ {- 1}) \psi (s) d \mu (s)
\]

是 \(\beta\)-可积的，且

\[
\begin{array}{l} \mathrm{I} = \iint f (s ^ {- 1} x) g (x) \chi (s ^ {- 1}) \psi (s) d \mu (s) d \beta (x) \\ = \int g (x) d \beta (x) \int f (s ^ {- 1} x) \chi (s ^ {- 1}) \psi (s) d \mu (s). \end{array}
\]

但 \(\psi \cdot \mu = \mu\)，因此

\[
\int f (s ^ {- 1} x) \chi (s ^ {- 1}) \psi (s) d \mu (s) = (\mu * ^ {\beta} f) (x)
\]

在局部 \(\beta\)-几乎处处成立。这表明函数

\[
x \mapsto g (x) \left(\left(\mu * ^ {\beta} f\right) (x)\right)
\]

是本质 \(\beta\)-可积的，且

\[
\mathrm{I} = \int g (x) \left(\left(\mu * ^ {\beta} f\right) (x)\right) d \beta (x).
\]

另一方面，引理 1 表明函数

\[
(s, x) \mapsto g (s x) f (x) \chi (s ^ {- 1}) \psi (s)
\]

对 \((\chi \cdot \mu) \otimes \beta\) 可积。因此函数 \((s, x) \mapsto g(s^{-1}x)f(x)\psi(s^{-1})\) 对 \(\breve{\mu} \otimes \beta\) 可积，且

\[
\begin{array}{l} \mathrm{I} = \iint g (s ^ {- 1} x) f (x) \psi (s ^ {- 1}) d \mu^ {\vee} (s) d \beta (x) \\ = \int f (x) d \beta (x) \int g (s ^ {- 1} x) \psi (s ^ {- 1}) d \mu^ {\vee} (s). \end{array}
\]

但 \(\check{\psi} \cdot \check{\mu} = \check{\mu}\)，于是 \(\int g(s^{-1}x)\psi(s^{-1})d\check{\mu}(s) = ((\chi \cdot \check{\mu}) *^{\beta}g)(x)\) 在局部 \(\beta\)-几乎处处成立。这表明函数

\[
x \mapsto f (x) \left(\left((\chi \cdot \stackrel {\vee} {\mu}) * ^ {\beta} g\right) (x)\right)
\]

是本质 \(\beta\)-可积的，且

\[
\mathrm{I} = \int f (x) \left(\left(\left(\chi \cdot \stackrel {\vee} {\mu}\right) * ^ {\beta} g\right) (x)\right) d \beta (x).
\]

这就证明了命题。

例. --- 4) 可取 \(f \in \mathcal{C}(\mathrm{X})\)，\(g \in \mathcal{K}(\mathrm{X})\) 与 \(\mu \in \mathcal{C}'(\mathrm{G})\)（取 \(\psi = 1\)）。

\begin{enumerate}
\def\labelenumi{\arabic{enumi})}
\setcounter{enumi}{4}
\item
  若 \(\mathbf{G}\) 真作用于 \(\mathbf{X}\)，则可取 \(f \in \mathcal{K}(\mathbf{X})\)，\(g \in \mathcal{K}(\mathbf{X})\) 与 \(\mu \in \mathcal{M}(\mathbf{G})\)（取 \(\psi = 1\)）。
\item
  可取 \(f \in \mathrm{L}^p(\mathrm{X},\beta)\)，\(g \in \mathrm{L}^q(\mathrm{X},\beta)\) 与 \(\mu \in \mathcal{M}^\rho(\mathrm{G})\)，其中 \(1 \leqslant p < +\infty\)，\(\frac{1}{p} + \frac{1}{q} = 1\)，\(\rho = \chi^{-1/q}\)。条件 (i) 与 (ii) 由命题 5 与命题 6 满足。我们来证明 (iii)。我们已经看到，\(\mu\) 载于一个集合 \(S\)，它是一列紧集的并。取 \(\psi\) 为 \(S\) 的特征函数。函数 \(h\) 是 \((\mu \otimes \beta)\)-可测的：因为函数 \((s,x) \mapsto g(x)\chi(s^{-1})\psi(s)\) 如此，且由引理 1，函数 \((s,x) \mapsto f(s^{-1}x)\) 也如此。此外，由于 \(g\) 在一列 \(\beta\)-可积集合的并之外为零，\(h\) 在一列 \((\mu \otimes \beta)\)-可积集合的并之外为零。于是我们有（第五章 §8, No.~3 命题 7）：
\end{enumerate}

\[
\begin{array}{l} \mathrm{J} = \iint^ {*} | g (x) f (s ^ {- 1} x) | \chi (s ^ {- 1}) \psi (s) d | \mu | (s) d \beta (x) \\ = \int^ {*} | g (x) | d \beta (x) \int^ {*} | f (s ^ {- 1} x) | \chi (s ^ {- 1}) \psi (s) d | \mu | (s). \end{array}\tag{9}
\]

而由于 \(g\)（相应地，\(\psi\)）在一列可积集合的并之外为零，(9) 式右端的上积分等于本质上积分（第五章 §1, No.~2 命题 7）。现在（第五章 §5, No.~3 命题 3）

\[
\int^ {\bullet} | f (s ^ {- 1} x) | \chi (s ^ {- 1}) \psi (s) d | \mu | (s) = \int^ {\bullet} | f (s ^ {- 1} x) | \chi (s ^ {- 1}) d | \mu | (s)
\]

由于 \(\mu = \psi \cdot \mu\)。由命题 6，最后这个积分有限，且等于 \((|\mu| *^{\beta} |f|)(x)\)（在局部 \(\beta\)-几乎处处）。因此

\[
J = \int^ {\bullet} | g (x) | \left(| \mu | * ^ {\beta} | f |\right) (x) d \beta (x),
\]

而 J 有限，因为 \(g \in \mathbf{L}^q\) 且 \(|\mu| *^\beta |f| \in \mathbf{L}^p\)（命题 6）。因此 \(h\) 是 \((\mu \otimes \beta)\)-可积的。

这时公式 (8) 表示自同态 \(g \mapsto (\chi \cdot \breve{\mu}) * g\)（\(\mathrm{L}^q(\mathrm{X},\beta)\) 的）对 \(\mu \in \mathcal{M}^\rho(\mathrm{G})\) 是自同态 \(f \mapsto \mu * f\)（\(\mathrm{L}^p(\mathrm{X},\beta)\) 的）的转置。

